% Generated draft: requires mathematical and source review.
\documentclass[11pt]{article}
\usepackage[margin=25mm]{geometry}
\usepackage{fontspec}
\setmainfont{Latin Modern Roman}
\usepackage{amsmath,amssymb,mathtools,unicode-math}
\setmathfont{Latin Modern Math}
\usepackage{enumitem,needspace,xurl,hyperref}
\setlength{\parindent}{0pt}
\setlength{\parskip}{5pt}
\setlength{\emergencystretch}{3em}
\newcommand{\N}{\mathbb{N}}
\newcommand{\Z}{\mathbb{Z}}
\newcommand{\Q}{\mathbb{Q}}
\newcommand{\R}{\mathbb{R}}
\newcommand{\C}{\mathbb{C}}
\newcommand{\F}{\mathbb{F}}
\newcommand{\A}{\mathbb{A}}
\newcommand{\PP}{\mathbb{P}}
\newcommand{\E}{\mathbb{E}}
\newcommand{\eps}{\varepsilon}
\newcommand{\dd}{\,\mathrm{d}}
\begin{document}
\section*{UnsolvedMath ID: 800715}
\textbf{Source problem code:} catalogue entry
\subsection{Definitions and mathematical statement}
Let \(F\) be a field, and let \(A\) and \(B\) be division \(F\)-algebras, as in the supplied statement. Under the usual associative, unital interpretation of an \(F\)-algebra, the tensor product \(T=A\otimes_F B\) has multiplication determined by
\[
(a\otimes b)(a'\otimes b')=(aa')\otimes(bb').
\]
Here a division algebra is an algebra in which every nonzero element has a multiplicative inverse. Say that \(T\) has a zero divisor if
\[
\exists\,x,y\in T\setminus\{0\}\quad xy=0.
\]
The problem asks for a structural criterion on \(A\) and \(B\) characterizing exactly when this condition holds.

The supplied record does not specify what counts as a structural criterion, nor whether the division algebras are required to be finite-dimensional over \(F\) or to have center exactly \(F\). These are unresolved qualifications of the intended problem. The source title places the question in the context of central simple algebras, but the extracted statement does not explicitly impose those hypotheses.
\subsection{Short English statement}
Find a structural characterization of when the tensor product of two division algebras over the same field contains nonzero elements whose product is zero. The intended dimensional and centrality hypotheses, and the meaning of a structural criterion, require clarification.
\subsection{Sources}
\begin{itemize}
\item Local UnsolvedMath registry: permanent ID \(800715\), source code \texttt{catalogue entry}, title \emph{Open Problems on Central Simple Algebras --- Problem 29}. The supplied record labels the problem open; this status is not independently verified.
\item Auel, Brussel, Garibaldi, and Vishne, \emph{Open Problems on Central Simple Algebras} (2010), Problem 29. Supplied source URL: \url{https://arxiv.org/abs/1006.3304}. The registry attributes the problem to these authors and gives proposed year 2010.
\item Supplied provenance reports access on 2026-10-16 and extraction by \texttt{arxiv\_latex}. This is supplied provenance, not a claim of independent source consultation. The record was created and updated on 2026-10-16; its upstream revision is \texttt{d371ee917ac9647fbb06cd6fc529deec6f3dda82}.
\end{itemize}
\end{document}
